% STATUS: DRAFT
\chapter{Haag--Kastler axioms}\label{ch:haagkastler}

\begin{physmotiv}
Quantum field theory is, at its core, the statement that observables can be \emph{localized}: associated with regions of spacetime, with those of spacelike separated regions commuting, and with Poincar\'e transformations moving regions around. The Haag--Kastler framework keeps exactly this and nothing more. Fields, Lagrangians, and even particle interpretations are derived or optional. It is the continuum, relativistic version of the quasi-local algebra of \cref{ch:quasilocal}, and it is the setting in which all statements about type III algebras of regions are made.
\end{physmotiv}

\section{The axioms}

A \emph{local net} assigns to each open bounded region $O\subset\R^{d}$ (typically double cones) a von Neumann algebra $\A(O)$ on a Hilbert space $\HH$ such that:
\begin{enumerate}[label=(HK\arabic*)]
\item \textbf{Isotony.} $O_1\subset O_2\Rightarrow\A(O_1)\subset\A(O_2)$.
\item \textbf{Locality (Einstein causality).} If $O_1$ and $O_2$ are spacelike separated, $[\A(O_1),\A(O_2)]=0$; equivalently $\A(O')\subset\A(O)'$ with $O'$ the causal complement. (For fermions, graded commutation.)
\item \textbf{Covariance.} There is a strongly continuous unitary representation $U$ of the Poincar\'e group with $U(g)\A(O)U(g)\ad=\A(gO)$.
\item \textbf{Spectrum condition.} The joint spectrum of the translations $U(a)=\ee^{\ii P\cdot a}$ lies in the closed forward light cone.
\item \textbf{Vacuum.} There is a unit vector $\Omega$, unique up to phase, invariant under $U$, and (irreducibility) cyclic for the algebra generated by all $\A(O)$.
\end{enumerate}
Optional structural properties that hold in many models: \emph{Haag duality} $\A(O')=\A(O)'$ (for wedges, or double cones in suitable models), \emph{additivity} (a region covered by smaller ones is generated by them), and the \emph{time-slice property} (the algebra of a neighbourhood of a Cauchy surface is the algebra of the whole domain of dependence).

The \emph{quasi-local algebra} $\A=\overline{\bigcup_O\A(O)}^{\norm\cdot}$ is a \cstar-algebra, the vacuum state $\omega_0=\inner\Omega{\,\cdot\,\Omega}$ its Poincar\'e-invariant state, and $(\HH,\pi_0,\Omega)$ the GNS data of $\omega_0$ (\cref{ch:gns}). Other representations (thermal states, charged sectors) arise from other states and are generally inequivalent (\cref{ch:inequiv}).

\begin{keyidea}
A QFT is a net of von Neumann algebras $O\mapsto\A(O)$ with isotony, locality, Poincar\'e covariance, positive energy, and a unique invariant vacuum. All dynamical content lives in how the algebras sit relative to each other.
\end{keyidea}

\section{Free-field example}

For the free scalar field of mass $m\ge0$, let $\mathcal S$ be the real symplectic space of Cauchy data and $\mathcal S_O$ the data supported in (the domain of dependence of) $O$. Define
\begin{equation}
\A(O)=\{W(f):f\in\mathcal S_O\}''\subset\Bnd{\mathcal F},
\end{equation}
with $W(f)=\ee^{\ii\phi(f)}$ the Weyl operators on Fock space (\cref{ch:weyl,ch:gns}). Isotony is obvious, locality follows from the Weyl relations because the symplectic form vanishes for spacelike supports, covariance from Poincar\'e invariance of the symplectic form and of the one-particle structure, and the spectrum condition from positivity of the one-particle energy. The Fock vacuum is the vacuum. This is the algebraic version of canonical quantization, defined without reference to a Hamiltonian.

\section{Relation to Wightman's framework}

In the Wightman approach one has operator-valued distributions $\phi(f)$ on a dense domain, with the same covariance and spectrum conditions, locality $[\phi(f),\phi(g)]=0$ for spacelike supports, and a cyclic vacuum. From such fields one \emph{builds} nets: $\A(O)$ is the von Neumann algebra generated by bounded functions (e.g.\ spectral projections or exponentials) of the self-adjoint $\phi(f)$ with $\mathrm{supp}f\subset O$, \emph{provided} the commutation of fields at spacelike separation holds in the strong sense (commutation of the spectral projections of the self-adjoint extensions), a technical point that is automatic for free fields and is an assumption in general. Conversely, from a net one cannot generally recover fields (these need additional smoothness), but the net carries all the physical information: scattering theory (Haag--Ruelle), the particle spectrum, and charges (Doplicher--Haag--Roberts, \cref{ch:dhr}).

\begin{whatwrong}
\begin{itemize}
\item The net is a statement about the \emph{vacuum representation}; ``the Hilbert space of the theory'' still means this representation. Haag's theorem (\cref{ch:why_algebras}) says that interacting vacua are not unitarily related to free ones, but \emph{nets} of interacting theories are fine.
\item Gravity does not fit: there is no Poincar\'e group acting on a fixed background and no local algebra of diffeomorphism-invariant observables (\cref{ch:gravity_constraints}). Parts VII--VIII show how the algebra of a region must be redefined.
\item Constructing interacting nets in $d=4$ is an open problem; they exist in $d=2$ and for conformal theories (\cref{ch:typeIII_local}), and free and many integrable models are covered.
\end{itemize}
\end{whatwrong}

\section*{Exercises}
\begin{exercise}
For the free field show that $[\A(O_1),\A(O_2)]=0$ for spacelike $O_1,O_2$ using the Weyl relation and the vanishing of the commutator function outside the light cone.
\end{exercise}
\begin{exercise}
Explain why an algebra $\A(O)$ cannot be type I (i.e.\ $\Bnd{\HH_O}$ for a tensor factor) if $\A(O')=\A(O)'$ and the vacuum is entangled across $\partial O$; anticipate \cref{ch:typeIII_local}.
\end{exercise}
\begin{exercise}
In $1+1$ dimensions with a massless free chiral boson, write $\A(I)$ for an interval $I$ of the light ray in terms of $W(f)$, $\mathrm{supp}f\subset I$, and check that translations and dilations act geometrically.
\end{exercise}
